\section*{第 \ref{ch:g11:forces-and-motion} 章——力与运动}

\begin{solution}{exo:g11:forces-and-motion:1}
(a) 速度不变：\omterm{def:g11:forces-and-motion:netforce}{合力}为零。(b) 速率改变：\omterm{def:g11:forces-and-motion:netforce}{合力}不为零，
与运动方向相反。(c) 方向改变（速度是矢量）：
\omterm{def:g11:forces-and-motion:netforce}{合力}不为零，指向弯道内侧。(d) 速度不变：\omterm{def:g11:forces-and-motion:netforce}{合力}为零——
\omterm{def:g10:universal-gravitation:weight}{重力}与阻力\omterm{def:g10:inertia:compensate}{相互抵消}。
\end{solution}

\begin{solution}{exo:g11:forces-and-motion:2}
$d_1/d_2 = 1.0/3.0$，且 $d_1 + d_2 = \qty{0.80}{m}$：
$d_1 = \qty{0.20}{m}$，为距 \qty{3.0}{kg} 小球的距离。两球质量相等时，
位于中点，距两球各 \qty{0.40}{m}。
\end{solution}

\begin{solution}{exo:g11:forces-and-motion:3}
$F\Delta t$ 相同，质量为三倍：装载后的 $\Delta v$ 为空车的三分之一。
若要与空车的 $\Delta v$ 相同，则推力需为三倍（或
作用时间需为三倍）。
\end{solution}

\begin{solution}{exo:g11:forces-and-motion:4}
$\Delta \vect v$：大小为 \qty{2}{m/s}，方向\emph{向后}（由
\num{10} 变为 \qty{8}{m/s}）。\omterm{def:g11:forces-and-motion:netforce}{合力}的方向与之相同：
与 $\vect v$ 相反——情形 2，减速。
\end{solution}

\begin{solution}{exo:g11:forces-and-motion:5}
(a) \omterm{def:g11:forces-and-motion:friction}{静摩擦力}，与你的推力方向相反（尚未滑动）。(b) \omterm{def:g11:forces-and-motion:friction}{动摩擦力}，
与滑动方向相反。(c) \omterm{def:g11:forces-and-motion:friction}{静摩擦力}，沿斜面\emph{向上}——
与\omterm{def:g10:universal-gravitation:weight}{重力}试图引起的下滑方向相反。
\end{solution}

\begin{solution}{exo:g11:forces-and-motion:6}
$\Delta v \propto F\Delta t/m$：(a) \qty{8.0}{m/s}；(b) \qty{2.0}{m/s}；
(c) \qty{2.0}{m/s}；(d) $4.0 \times \frac{3 \times 2}{6} =
\qty{4.0}{m/s}$——三次变化\omterm{def:g10:inertia:compensate}{相互抵消}。
\end{solution}

\begin{solution}{exo:g11:forces-and-motion:7}
$v' = \sqrt{8.0^2 + 6.0^2} = \qty{10.0}{m/s}$，此时 $\tan\theta =
6.0/8.0$，即 $\theta \approx 37^\circ$，为东偏北。速度的合成
应按\emph{矢量}进行：相互垂直的分量由毕达哥拉斯定理合成，
而不是直接把大小相加。
\end{solution}

\begin{solution}{exo:g11:forces-and-motion:8}
事先固定的量：小球的质量及其 $\Delta v$（入射速率
减为零），因而乘积 $F\Delta t$ 固定。使 $\Delta t$ 由
\qty{0.02}{s} 增至 \qty{0.08}{s}，平均力除以 $4$。
\end{solution}

\begin{solution}{exo:g11:forces-and-motion:9}
鞋受到的是跑道对鞋的\omterm{def:g11:forces-and-motion:friction}{静摩擦力}：脚向后蹬地，
地面与脚之间的相互\omterm{def:g10:inertia:contact}{接触力}推动运动员向前。钉鞋提高
最大\omterm{def:g11:forces-and-motion:friction}{静摩擦力}（全力时不打滑）。在无摩擦的冰面上
不存在水平方向的外力，因此
\omterm{def:g11:forces-and-motion:com}{质心}速度无法改变——既无法起跑，也无法行走。
\end{solution}

\begin{solution}{exo:g11:forces-and-motion:10}
滚动车轮的接触点不滑动：受\omterm{def:g11:forces-and-motion:friction}{静摩擦力}，
该力大于打滑时的\omterm{def:g11:forces-and-motion:friction}{动摩擦力}，
（见\cref{def:g11:forces-and-motion:friction}）——停车距离更短。
打滑的轮胎所受\omterm{def:g11:forces-and-motion:friction}{动摩擦力}与滑动方向相反，无论
车轮朝向如何。转动车轮无济于事，转向能力丧失。
\end{solution}

\begin{solution}{exo:g11:forces-and-motion:11}
钢丝的\omterm{def:g10:inertia:inventory}{张力}指向运动员。松手后
小球沿切线方向做直线运动飞出，速率
与原来相同（惯性）。月球的“钢丝”是地球的
引力；若关闭引力，月球将沿切线做\omterm{def:g10:inertia:uniform}{匀速直线运动}，
永远如此。
\end{solution}

\begin{solution}{exo:g11:forces-and-motion:12}
从顶端算起：$x = \dfrac{1.0 \times 0.60 + 2.0 \times 1.20}{3.0}
= \qty{1.00}{m}$。两部分\omterm{def:g11:forces-and-motion:com}{质心}的距离分别为 \qty{0.40}{m}
（杆）和 \qty{0.20}{m}（头部），比值 $2 : 1$ 恰为
质量比 $1.0 : 2.0$ 的倒数，正如\cref{def:g11:forces-and-motion:com}所要求。
%
\end{solution}

\begin{solution}{exo:g11:forces-and-motion:13}
\emph{1.} $F$ 与 $\Delta v$ 相同，$m$ 加倍，货车所需
时间为两倍。\emph{2.} \omterm{def:g10:relative-motion:speed}{平均速率}相同，时间为两倍时，
距离也为两倍。\emph{3.} $\Delta v$ 加倍，而 $F$ 与 $m$ 不变：所需
时间为两倍；但\omterm{def:g10:relative-motion:speed}{平均速率}也加倍，因此距离
变为原来的 $4$ 倍。速率加倍时，
制动距离增至四倍。
\end{solution}

\begin{solution}{exo:g11:forces-and-motion:14}
\emph{1.} $T = 27.3 \times \qty{86400}{s} \approx \qty{2.36e6}{s}$，故
$v = \dfrac{2\pi \times \num{3.84e8}}{\num{2.36e6}} \approx
\qty{1.0e3}{m/s} \approx \qty{1}{km/s}$。\emph{2.} 指向地球
（向心）；即地球的引力。\emph{3.} 月球\emph{正在}
下落：其速度变化始终指向地球。
其切向速率只保证它不断错过地球——它
是\emph{绕着}地球下落，而非落入地球。
\end{solution}

\begin{solution}{exo:g11:forces-and-motion:15}
$m$ 与 $\Delta v$ 相同，故 $F\Delta t$ 固定：在泡沫上，$\Delta t$
延长为原来的五十倍，平均力便缩小为原来的五十分之一。
屈膝可使着地时的 $\Delta t$ 由一次猛击
延长为较长的缓冲过程。从力学上讲，防撞垫
实际上是由额外的数毫秒“堆”出来的。
\end{solution}

\begin{solution}{pb:g11:forces-and-motion:1}
\textbf{1.} $50/3.6 \approx \qty{13.9}{m/s}$（记为 \qty{14}{m/s}）。
\textbf{2.} $\Delta v \approx \qty{14}{m/s}$：乘员碰前以
\qty{14}{m/s} 运动，碰后为零——两端均由碰撞确定，
因此没有任何装置能改变 $\Delta v$。
\textbf{3.} $\Delta v \propto F\Delta t/m$
由\cref{thm:g11:forces-and-motion:secondlaw}可知，在 $\Delta v$ 与 $m$
固定时，$F\Delta t$ 也须固定。设计者唯一能\omterm{def:g11:lenses-and-eye:eye}{调节}的量是
$\Delta t$。
\textbf{4.} $0.150/0.050 = 3$：吸能区使平均
力除以 $3$。
\textbf{5.} 向后，与运动方向相反——速度由
\qty{14}{m/s} 减小为零。
\textbf{6.} 折叠前端的作用，是让停止过程
变得更长：会变形的金属能换来数毫秒。
\textbf{7.} $\Delta v$ 与 $m$ 相同，$\Delta t$ 缩短为原来的三分之一：
老式假人承受的力是原来的三倍。
\textbf{8.} \omterm{def:g10:relative-motion:speed}{平均速率} $\approx \qty{6.9}{m/s}$；距离 $\approx
6.9 \times 0.150 \approx \qty{1.0}{m}$——恰好约为
一米左右的可变形车身，正是真实吸能区所提供的。
\textbf{9.} 一个 \qty{10}{m} 长的车头既不能驾驶，也无法停放，且其
自身质量会加剧每次碰撞。乘员舱决不能
变形：它须保留生存空间。
\textbf{10.} “\dots 因此你必须选择 $\Delta t$。”
\textbf{11.} $0.150/0.005 = 30$：挡风玻璃造成的停止
猛烈程度是系安全带者的三十倍。
\textbf{12.} 略有伸展的安全带会使乘员自身的
$\Delta t$ 再略为延长（并使力分散）；钢制束带
则会把车身最短的停止过程强加于最柔弱的乘员。
\textbf{13.} $0.080/0.010 = 8$：安全气囊使头部所受的力除以
$8$。
\textbf{14.} $\Delta v$ 与 $\Delta t$ 相同，力与
质量成正比，成人所受的力大 $70/20 = 3.5$ 倍。
\textbf{15.} 所需的速度变化：$14/0.150 \approx \qty{93}{m/s}$ 每
秒——约为\omterm{def:g10:universal-gravitation:weight}{重力}所产生量的 $9.5$ 倍
（\qty{9.8}{m/s} 每秒）。双臂须施加的拉力约为
婴儿\omterm{def:g10:universal-gravitation:weight}{重力}的 $9.5$ 倍，好比托举一个 \qty{95}{kg} 的重物。无人手臂能做到；婴儿
则向前飞出。
\textbf{16.} 成人所用的安全带将力传递至髋部和胸部；若用于
儿童，则会使力作用于腹部和颈部。儿童座椅施加的
（较小）力作用于强壮的身体部位，其垫层使
$\Delta t$ 进一步延长。
\textbf{17.} 若乘员舱发生变形，乘员将失去
生存空间，并被结构本身撞击。任何东西都不应直接撞击
刚性的乘员舱：乘员的停止过程由安全带与
安全气囊承担（较长的 $\Delta t$），而乘员舱则固定住会变形的两端。
\textbf{18.} 由对称性，两车均在接触面处停止：每位
驾驶员在吸能区内的停止过程中经历 $\Delta v \approx \qty{14}{m/s}$ 的
速度变化——本质上是撞墙测试，而非“双倍碰撞”。
\textbf{19.} 两力大小相等、方向相反，但
$\Delta v \propto F\Delta t/m$：小车质量 $m$ 小，
所经历的速度变化却大。
\textbf{20.} 吸能区、可伸展的安全带、安全气囊、儿童座椅
的垫层——它们都在延长同一个量：$\Delta t$。
\textbf{21.} 碰撞固定了质量与速度变化，因此
力 $\times$ 时间这一乘积固定：轻柔的碰撞即是时间较长的碰撞——
延长时间，力便按相同比例减小。
\end{solution}
